\section{从 GPT-2 到 GPT-3：把直觉变成可检验 Scaling Hypothesis}

Mann 把 2015 ImageNet 和随后 GPT-2 的出现视为关键 field signal：长期被认为只是 academic promise 的神经网络开始在真实任务上表现出可迁移能力。GPT-2 提示“在大规模互联网文本上做 next-token prediction 可能产生广泛能力”，GPT-3 则用多尺度实验检验参数、数据和 compute 增长是否稳定改善 loss 与 few-shot performance。

\subsection{Observation、hypothesis、pilot 与 frontier run}

\term{scaling law}（缩放定律）是模型性能或 loss 与 compute、dataset size、parameter count 之间的经验规律。它的工程价值不是事后画一条漂亮线，而是在最大训练前用较小 run 估计收益、比较 recipe、配置 data/model mix，并明确何时新证据表明 regime 已改变。

\lecturefigure{02-scaling-hypothesis.png}{Scaling law 把 field signal 转成 hypothesis、pilot runs 和可预测的 frontier commitment。}{本地字幕 01:00--05:40；GPT-2/GPT-3 与 scaling laws 论文。}

\begin{knowledgebox}{读图：最大 run 之前必须有预测}
Field signal 提供方向，hypothesis 明确“哪些变量增加会怎样改变 loss”，pilot runs 要跨多个尺度并保持可比 recipe，frontier run 则在预算批准前写出 expected curve、confidence 和 stop condition。若团队只有一个大 run，没有 holdout prediction，就很难判断成功来自 scaling、data cleanup、architecture change 还是选择性解释。
\end{knowledgebox}

\teachervoice{Mann 强调自己没有传统 PhD 路径，却通过 data engineering、分析与 architecture experiment 参与 GPT-3。课堂提示：frontier research 的关键贡献不只来自提出新公式，也来自构造可信数据、实验和测量系统。}

\subsection{为什么当时很多人不信}

怀疑并不荒谬：历史技术常经历 sigmoid，已有模型可能成本高、部署困难，BERT/T5 时代的经验又强调 task-specific fine-tuning。问题在于把“当前 recipe 和 hardware 下的经济性”误当成“永远的能力上限”。更好的争论方式是要求跨尺度预测、residual analysis 与外推边界，而不是用一次昂贵 demo 或一次失败否定整个假设。

\begin{warningbox}{Scaling law 不是“只要加 compute 就一定成功”}
Architecture、optimizer、token mixture、data contamination、context length、precision、parallelism 和 post-training 都可能改变曲线。Power law 只描述已观测区间的经验关系；跨 modality、跨 recipe 或跨多个数量级外推时，应扩大 uncertainty，并准备发现新的 bottleneck。
\end{warningbox}

\subsection{本章小结}

Scaling hypothesis 的科学性来自可预测和可证伪。Field signal 触发方向，多尺度 pilot 提供曲线，frontier run 检验外推；任何结果都必须与预注册预测和 residual 对照。

